\section{Related Work}

\paragraph{LoRA adaptation and parameter-efficient fine-tuning security.}
LoRA reduces fine-tuning cost by freezing the base model and representing task updates through low-rank matrices injected into selected transformer projections \citep{hu2022lora}.
QLoRA combines a frozen 4-bit quantized base model with trainable low-rank adapters, showing that adapter fine-tuning and low-precision deployment can be coupled \citep{dettmers2023qlora}.
Recent parameter-efficient fine-tuning work makes the adapter object more fine-grained or lower-bit: GraLoRA allocates rank at a finer granularity, and LowRA targets LoRA fine-tuning under 2-bit constraints \citep{jung2025gralora,zhou2025lowra}.
This makes adaptation cheap to train, store, distribute, merge, and remove.
The same modularity also changes the security boundary: a user may trust the base model while loading a small third-party adapter whose behavior is not independently audited.
Adapter-level attacks are therefore a natural supply-chain risk.
Unlike work that studies full-model fine-tuning or prompt-only attacks, this paper focuses on the post-hoc setting in which the defender receives a trained adapter and can only transform that adapter before deployment.

\paragraph{Large language model backdoors and sleeper behavior.}
Backdoor and sleeper-agent studies show several routes by which language models can learn conditional behavior: weight poisoning modifies pretrained models \citep{kurita2020weight}, BadPrompt attacks continuous prompts \citep{cai2022badprompt}, BadEdit implants backdoors through model editing \citep{li2024badedit}, and sleeper-agent training studies deceptive LLM behavior that can persist through safety training \citep{hubinger2024sleeper}.
Recent surveys summarize the broader LLM backdoor attack and defense landscape \citep{zhang2024llmbackdoorsurvey,zhou2025backdoorsurvey}.
BackdoorLLM provides a benchmark and evaluation pipeline for generative LLM backdoors \citep{li2025backdoorllm}.
JailbreakLoRA studies malicious LoRA adapters from sharing platforms \citep{wei2026jailbreaklora}, while \citet{salles2025lorausers} show that a few spurious tokens injected during parameter-efficient fine-tuning can manipulate test-time behavior.
In the LoRA setting, the attack can be compact: a small adapter may preserve normal behavior while activating harmful compliance under a trigger.
This makes attack-success-only evaluation insufficient.
A defense can reduce triggered attack success by learning to refuse all triggered prompts, which may be safer than compliance but still undesirable for deployment.
Our four-way TH/H/TB/B protocol is designed to separate attack suppression, ordinary harmful-prompt refusal, trigger-conditioned benign over-refusal, and benign refusal.

\paragraph{Backdoor defenses and safety tuning.}
Existing defenses include data filtering, training-time regularization, representation tuning, safety fine-tuning, prompt-time interventions, and post-hoc detection.
Safety-tuning work such as instruction tuning, reinforcement learning from human feedback, harmlessness training, and constitutional AI provides the broader alignment background \citep{ouyang2022training,bai2022helpful,bai2022constitutional}.
Representation-space interventions and refusal-direction analyses show that safety behaviors can sometimes be controlled through low-dimensional activation directions \citep{zou2023representation,arditi2024refusal}.
Recent post-training defenses further explore internal-consistency regularization, post-fine-tuning perturbations, LoRA safety pruning, and trigger-agnostic backdoor purification \citep{min2025crow,wang2025panacea,ao2025splora,li2026purifying}.
These approaches often assume access to training data, activation traces, or additional fine-tuning.
\sac{} instead acts directly on adapter weights after training.
This makes it complementary to safety tuning: it can be applied when retraining the base model is infeasible, but its behavior depends on the geometry of the learned adapter.
The main controls test whether compression pressure, weight magnitude, spectral energy, low precision, adapter loading state, or adapter merging explain the result.
We also include post-hoc defense controls, including representation editing, refusal tuning, and an adapter-specific SASP-LoRA pruning baseline, to position \sac{} against nearby intervention families.

\paragraph{Compression, pruning, and quantization.}
Compression methods such as pruning, quantization, and low-rank approximation are usually evaluated through efficiency and utility.
Deep Compression combines pruning, quantization, and coding for smaller neural networks \citep{han2016deepcompression}.
LLM.int8 and SmoothQuant study low-precision LLM inference \citep{dettmers2022llmint8,xiao2023smoothquant}, GPTQ provides post-training quantization for generative transformers \citep{frantar2023gptq}, and SparseGPT/Wanda study one-shot LLM pruning \citep{frantar2023sparsegpt,sun2024wanda}.
QLoRA and LowRA connect low-precision operation with adapter fine-tuning \citep{dettmers2023qlora,zhou2025lowra}, while SVD-LLM V2 studies singular-value truncation for LLM compression \citep{wang2025svdllmv2}.
For backdoored adapters, utility alone is not enough.
Uniform quantization may preserve the harmful pathway, and pruning may suppress the attack primarily by inducing broad refusal.
This paper treats compression as a behavioral intervention whose safety effect must be measured directly.
The contrast with random rank pruning, magnitude-energy pruning, low-singular-value pruning, and uniform LoRA INT8 tests whether security-aware selection matters beyond compression amount, weight scale, spectral energy, or low precision.

\paragraph{Editing and causal localization.}
Model editing methods aim to localize and alter internal mechanisms that support a target behavior \citep{meng2022rome,meng2023memit}.
Representation-engineering and refusal-direction work pursue a related goal through activation-space interventions \citep{zou2023representation,arditi2024refusal}.
Mechanistic and causal analyses similarly study how internal components contribute to model outputs \citep{elhage2021transformer,huben2024sparseautoencoders}.
\sac{} is related in that it uses counterfactual interventions to decide which components matter.
We use behaviorally grounded component scores as a practical deployment criterion, then validate the resulting adapter with held-out behavioral evaluation.
Recent backdoor-purification work also points to structured trigger--behavior associations inside model components, but targets a broader full-model or adapter-only purification setting rather than selective adapter compression \citep{li2026purifying}.
The Llama results show why this validation step matters: the same class of intervention can improve TH while revealing a higher-cost utility/refusal frontier.

\paragraph{Positioning.}
The closest conceptual view is post-training adapter editing under a compression budget.
The contribution is not that pruning or quantization alone removes backdoors; the experiments show the opposite for several controls.
Rather, the contribution is a security-aware selection rule and evaluation protocol for deciding which adapter directions should be removed, attenuated, or stored differently.
